\documentclass[11pt,a4paper]{article}

\usepackage[fontset=fandol]{ctex}
\usepackage{amsmath,amssymb}
\usepackage{graphicx}
\usepackage[margin=2.35cm]{geometry}
\usepackage[most]{tcolorbox}
\usepackage{booktabs,longtable,tabularx}
\usepackage{float}
\usepackage{xcolor}
\usepackage{hyperref}
\usepackage{fancyhdr}
\usepackage{etoolbox}
\usepackage{listings}

\definecolor{ink}{HTML}{25313D}
\definecolor{muted}{HTML}{66727E}
\definecolor{directBg}{HTML}{EDF5F1}
\definecolor{directFrame}{HTML}{6B9B82}
\definecolor{claimBg}{HTML}{EEF3F8}
\definecolor{claimFrame}{HTML}{527A9E}
\definecolor{reconBg}{HTML}{FBF5E8}
\definecolor{reconFrame}{HTML}{A77A2D}
\definecolor{uncertainBg}{HTML}{FAEEEE}
\definecolor{uncertainFrame}{HTML}{A65B5B}

\hypersetup{colorlinks=true,linkcolor=claimFrame,urlcolor=claimFrame,citecolor=claimFrame}
\setlength{\parindent}{0pt}
\setlength{\parskip}{0.65em}
\pagestyle{fancy}
\fancyhf{}
\fancyhead[L]{\small\textcolor{muted}{\reportshorttitle}}
\fancyhead[R]{\small\textcolor{muted}{Slides Reconstruction}}
\fancyfoot[C]{\thepage}
\renewcommand{\headrulewidth}{0.3pt}

\tcbset{evidencebox/.style={
  enhanced,breakable,boxrule=0.75pt,arc=1.5pt,
  left=3mm,right=3mm,top=2mm,bottom=2mm,
  fonttitle=\bfseries,coltitle=white
}}
\newtcolorbox{sourcebox}[1]{evidencebox,colback=directBg,colframe=directFrame,colbacktitle=directFrame,title={E0 直接证据：#1}}
\newtcolorbox{claimbox}[1]{evidencebox,colback=claimBg,colframe=claimFrame,colbacktitle=claimFrame,title={E1 跨页结论：#1}}
\newtcolorbox{reconstructbox}[1]{evidencebox,colback=reconBg,colframe=reconFrame,colbacktitle=reconFrame,title={E2 可能讲述：#1}}
\newtcolorbox{uncertaintybox}[1]{evidencebox,colback=uncertainBg,colframe=uncertainFrame,colbacktitle=uncertainFrame,title={不确定性：#1}}

\lstset{
  basicstyle=\ttfamily\small,
  breaklines=true,
  frame=single,
  numbers=left,
  numberstyle=\tiny\color{muted},
  captionpos=b
}

% Fill these fields before compiling.
\newcommand{\reporttitle}{[课程或幻灯片标题]}
\newcommand{\reportshorttitle}{[短标题]}
\newcommand{\reportsubtitle}{基于幻灯片证据的课程重建笔记}
\newcommand{\sourceauthor}{[作者或机构，未知则写“未确认”]}
\newcommand{\sourcefiles}{[源文件名列表]}
\newcommand{\reconstructiondate}{\today}
\newcommand{\coverimagepath}{}

\begin{document}

\begin{titlepage}
\centering
\vspace*{1.0cm}
{\Huge\bfseries\color{ink}\reporttitle\par}
\vspace{0.5cm}
{\Large\color{muted}\reportsubtitle\par}
\vspace{0.9cm}

\ifdefempty{\coverimagepath}{}{
  \includegraphics[width=0.9\textwidth,height=0.48\textheight,keepaspectratio]{\coverimagepath}\par
  \vspace{0.7cm}
}

\vfill
\begin{tcolorbox}[width=0.92\textwidth,colback=black!2,colframe=black!45,arc=1.5pt]
\textbf{原作者/机构}：\sourceauthor\par
\textbf{源文件}：\sourcefiles\par
\textbf{重建日期}：\reconstructiondate\par
\textbf{说明}：本报告依据幻灯片重建最可能的课程论证，不声称恢复演讲者的逐字原话。
\end{tcolorbox}
\end{titlepage}

\section*{阅读说明}
\addcontentsline{toc}{section}{阅读说明}
本文使用四级证据标记：E0 为页面直接证据，E1 为跨页强推断，E2 为可能的现场讲述，E3 为外部背景。推断与外部补充不得被当作原作者的确定表述。

\tableofcontents
\newpage

% Required outline:
% \section{课程论证地图}
% \section{... reconstructed teaching sections ...}
% \section{逐页重建索引}
% \section{总结、边界与延伸}
%
% Use sourcebox, claimbox, reconstructbox, and uncertaintybox only when the
% distinction materially helps the reader. Keep figures outside boxes.
% Every figure caption must identify the source deck and slide number.

\section{课程论证地图}
[在此写入核心问题、最可能的中心主张、章节关系和关键不确定性。]

\section{逐页重建索引}
[在此覆盖每一页的功能、核心观点、证据等级、置信度和正文位置。]

\section{总结、边界与延伸}
[在此总结最可靠的结论、论证连接、缺失前提、歧义以及需要原始录像或作者确认的问题。]

\end{document}
